\documentclass[10pt]{article}
\usepackage{german,mathsym,amsfonts}
\newtheorem{definition}{Definition}[section]
\newtheorem{theorem}{Theorem}[section]
\newtheorem{satz}[theorem]{Satz}
\newtheorem{lemma}[theorem]{Lemma}
\newtheorem{anmerk}{Anmerkung}[section]
\newtheorem{bsp}{Beispiel}[section]
\newtheorem{cor}[theorem]{Corollary}
\newtheorem{alg}{Algorithm}[section]
\title{Zusammenfassung der "Ubung zu Entropie}
\author{Pawel Wocjan\footnote{Bei Fragen, Verbesserungsvorschl"agen, etc. 
email an mich: \texttt{wocjan@ira.uka.de}}}
\begin{document}
\pagestyle{plain}
\maketitle
\tableofcontents
\pagenumbering{arabic}
\pagestyle{headings}


\section{Kombinatorische Grundbegriffe}
\subsection{Stirling-Approximation}
\begin{definition}[Fakult"at]
  Die f"ur alle nat"urlichen Zahlen $n\in\N$ definierte Funktion $f(n)=n!$, 
  f"ur die gilt: $f(0)=1$ und f"ur alle $n$: $f(n+1)=(n+1)\cdot f(n)$, hei"st
  $n$-Fakult"at.
\end{definition}
Zur n"ahrungsweisen Berechnung von $n!$ f"ur sehr gro"se Zahlen $n$ benutzt
man die \emph{Stirlingsche Formel}.
\begin{satz}[Stirlingsche Formel]
Die Stirlingsche N"ahrung lautet
$$
  n! \approx \sqrt{2\pi n}  \cdot n^n \cdot e^{-n} \quad\mbox{bzw.}\quad
  \ln(n!)\approx \left(n+{1\over 2}\right)\ln n - n + {1\over 2}\ln(2\pi) \,.
$$
\end{satz}
%
Beweis: (aus \cite{ham}) Zur Herleitung der Stirlingschen Formel, beweisen wir 
zun"achst die folgenden Absch"atzungen:
\begin{eqnarray*}
  n! & <    & e^1 \sqrt{n}            \cdot n^n \cdot e^{-n} \\
  n! & >    & e^{{7\over 8}} \sqrt{n} \cdot n^n \cdot e^{-n}
\end{eqnarray*}
Mit einem Trick k"onnen wir sie dann benutzen, um die Stirlingsche Formel
herzuleiten.

Der Beweis geht "uber die geometrisch anschauliche Absch"atzung 
von $\int_1^n \ln x dx$ durch eine Obersumme $O$ und eine Untersumme $U$. 
Es gilt
$$
  \int_{1}^n \ln x dx = n\ln n - n + 1.
$$
Anhand des folgenden Bildes (selber "uberlegen) wird die Untersumme $U$ 
entwickelt. Wir approximieren den Fl"acheninhalt durch die Trapezregel:
$$
  \int_1^n \ln x\, dx \ge \sum_{i=1}^{n-1} {\ln i + \ln(i+1)\over 2}
$$
\begin{eqnarray*}
  U & = & {1\over 2}\ln n + {1\over 2}\ln 1 + \sum_{i=2}^{n-1}\ln i \\
    & = & \ln(n-1)! + {1\over 2}\ln n \\
    & = & \ln n! - {1\over 2}\ln n
\end{eqnarray*}
Wir haben also $\ln n! - {1\over 2}\ln n < n\ln n - n + 1$. Wir l"osen nach
$\ln n!$ auf und erhalten so
$$ \ln n! < 1+{1\over 2}\ln n + n\ln n - n \, . $$
Wir wenden die Exponentialfunktion auf diese Ungleichung und erhalten so die
\emph{obere Schranke}
$$
  n! < n^n e^{-n} \sqrt{n} e \le n! \, .
$$
Die Berechnung f"ur die Obersumme $O$ ist etwas aufwendiger. Das folgende
Bild macht sie anschaulich.
Wir unterscheiden drei F"alle.
\begin{itemize}
  \item Der Fl"acheninhalt "uber $[1,1.5]$ wird durch den des Dreiecks
    "uberboten, dessen Hypothenuse die Steigung $1$ hat. Dieses Dreieck
    hat den Fl"acheninhalt $O_1={1\over 8}$.
  \item "Uber den Intervallen $[i-{1\over 2}, i+{1\over 2}]$ f"ur 
    $i=2,\ldots,n-1$ bilden wir Trapeze, deren Mittelinie die
    L"ange $\ln i$ hat und deren Oberseite die Tangentensteigung der
    Logarithmuskurve als Steigung hat. Der Fl"acheninhalt ist
    H"ohe (hier die Intervall"ange) mal Mittelinie gleich $\ln i$. 
    Alle Trapeze haben also insgesamt den Fl"acheninhalt
    $O_2 = \sum_{i=2}^{n-1}\ln i = \ln(n-1)!$.
  \item "Uber $[n-{1\over 2},n]$ bilden wir das Rechteck, dessen Seitenl"ange
    $\ln n$ ist. Es hat den Fl"acheninhalt $O_3={1\over 2}\ln n$.
\end{itemize}
So ergibt sich insgesamt f"ur die Obersumme 
\begin{eqnarray*}
  O & = & O_1 + O_2 + O_3 \\
    & = & {1\over 8} + \ln n! - {1\over 2}\ln n
\end{eqnarray*}
und damit 
$$
  {1\over 8} + {1\over 2}\ln n + n\ln n < \ln n! \, ,
$$
wenn wir nach $\ln n!$ aufl"osen. Wir wenden die Exponentialfunktion auf diese
Ungleichung und erhalten so die \emph{untere Schranke}
$$
  n! > n^n e^{-n} \sqrt{n} e^{7/8} \, .
$$
Nach den Ungleichung ist $n!$ also $n^n e^{-n} \sqrt{n} C$, wobei $C$
zwischen $e^{{7\over 8}}$ und $e$ liegt. Wir beweisen noch, da"s
$$ \lim_{n\rightarrow\infty} C = \sqrt{2\pi} $$
gilt.

Um diesen Koeffizienten $C$ f"ur $n\rightarrow\infty$ zu berechnen,
wenden wir einen Trick an und betrachten dazu das Integral
$$ I_k = \int_0^{\pi/2} \cos^k x dx $$
f"ur $k\ge 0$. Die trigonometrische Identit"at $\cos^2 x = 1-\sin^2 x$ 
liefert f"ur $k\ge 2$
\begin{eqnarray*}
  I_k & = & \int_0^{\pi/2} \cos^{k-2} x (1-\sin^2 x) dx \\
      & = & I_{k-2} + \int_0^{\pi/2} \cos^{k-2} x (-\sin x) \sin x dx \\
      & = & I_{k_2} + 
            {1\over k-1} \int_0^{\pi/2} (\cos^{k-1} x)'\cdot\sin x dx \, .
\end{eqnarray*}
Partielle Integration ergibt nun
\begin{eqnarray*}
  I_k & = & I_{k-2} + {\cos^{k-1} x\over k-1} \sin x\vert_0^{\pi/2} -
            {1\over k-1} \int_0^{\pi/2} \cos^{k-1} x \cos x dx \\
      & = & I_{k-2} - {1\over k-1} I_k\,.
\end{eqnarray*}
Bei der Aufl"osung nach $I_k$ erhalten wir die Rekurrenzformel 
$$
  I_k = {k-1\over k} I_{k-2}
$$
f"ur $k\ge 2$. Man kann leicht zeigen, da"s
\begin{eqnarray*}
  I_1 & = & 1 \\
  I_0 & = & {\pi \over 2} 
\end{eqnarray*}
gilt.

Wir bemerken nun, da"s wegen $\cos x\le 1$ ($x\in [0,{\pi\over 2}]$) gilt
$$ I_{2k-1} > I_{2k} > I_{2k+1}\,. $$
Wir teilen diese Ungleichung durch $I_{2k-1}$ und erhalten 
$$ 1 > {I_{2k}\over I_{2k-1}} > {I_{2k+1}\over I_{2k-1}}\, . $$
Wir ersetzen sukzessive die Terme mit Hilfe der Rekurrenzformel und erhalten
$$ 1 > 2k \left[ {2k-1\over 2k} {2k-3\over 2k-2} \cdots {3\over 2}\right]^2
       {\pi\over 2}
     > {2k\over 2k+1}  $$
Wir k"onnen den Z"ahler in der eckigen Klammer gleich $(2k)!$ machen, indem
wir den Bruch mit dem dem Nenner erweitern. Wir verwenden
$$ 2k(2k-2)(2k-4)\cdots = 2^k k! $$
und erhalten
$$ 1 > \pi k \left[ {(2k)! \over 2^k\cdot (k!)\cdot 2^k\cdot (k!)} \right]^2 
     > {2k\over 2k+1} \, .$$
Wir ziehen nun aus jedem Glied die positive Quadratwurzel 
$$ 1 > \sqrt{\pi k}\left[{(2k)!\over (k!)^2\cdot 2^{2k}}\right]
     > \sqrt{1-{1\over 2k+1}} $$
Schlie"slich setzen wir die N"ahrung
$$ n! = n^n e^{-n} \sqrt{n} C $$
in diese Ungleichung f"ur jede Fakult"at ein und erhalten 
$$ 1 > \sqrt{\pi k}\left[{ (2k)^{2k} e^{-2k}\sqrt{2k} C\over
       2^{2k} k^k e^{-k} \sqrt{k} C k^k e^{-k} \sqrt{k} C} \right] 
     > \sqrt{1-{1\over 2k+1}} $$
Deshalb gilt f"ur $k\rightarrow\infty$
$$ C\rightarrow \sqrt{2\pi}. $$
Damit haben wir die Stirlingsche Formel bewiesen
$$
  n\approx\sqrt{2\pi n} \cdot n^n \cdot e^{-n} \, .
$$
Das Verh"altnis des Stirlingschen N"ahrungswertes zum wahren Wert ist selbst
f"ur kleine Werte bemerkenswert gut. Die N"ahrung ist so beschaffen, da"s sie
f"ur $n\rightarrow\infty$ \emph{relative} exakt ist, die Differenz jedoch 
\emph{nicht} Null erreicht.

Wenn dies ungewohnt erscheint, so betrachte man die beiden Funktionen
\begin{eqnarray*}
  f(n) & = & n + \sqrt{n} \\
  g(n) & = & n
\end{eqnarray*}
Der Limes des Verh"altnisses ist
$$ {f(n)\over g(n)}={n+\sqrt{n}\over n}=1+{1\over\sqrt{n}}\rightarrow 1\, .$$
Die Differenz betr"agt jedoch 
$$ f(n)-g(n)=\sqrt{n}\rightarrow\infty \, .$$
Die Approximation $n\approx\sqrt{2\pi n} \cdot n^n \cdot e^{-n}$ ist schon 
f"ur kleine Werte von $n$ genau, aber es gibt noch sch"arfere 
Absch"atzungen
$$
  \sqrt{2\pi n} \cdot n^n \cdot e^{-n} \cdot e^{\frac{1}{12n+1}} < n ! <
  \sqrt{2\pi n} \cdot n^n \cdot e^{-n} \cdot e^{\frac{1}{12n}}\, ,
  \label{scharfStirling}
$$
deren Beweis in \cite{krengel91} gefunden werden kann.

\subsection{Binomial- und Multinominalkoeffizienten}
\begin{definition}[Binominalkoeffizient]
  Die f"ur alle nat"urlichen Zahlen $n$ und $k$ ($0\le k\le n$) definierte 
  Funktion 
  $$
    {n\choose k} = {n!\over k!\, (n-k)!} 
  $$
  hei"st Binominalkoeffizient.
\end{definition}
Die Werte des Binominalkoeffizienten k"onnen sukzessiv aus dem 
\emph{Pascalschen Dreieck} ermittelt werden.
\begin{definition}[Multinominalkoeffizient]
  Die f"ur alle nat"urlichen Zahlen $n$ und alle $q$-Tupel 
  $(k_1,k_2,\ldots,k_q)$, f"ur die $\sum_{i=1}^q k_i=n$ gilt, erkl"arte
  Funktion
  $$
    {n\choose k_1, k_2,\ldots, k_q}={n!\over k_1!\, k_2! \cdots k_q!}
  $$
  hei"st Multinominalkoeffizient.
\end{definition}
%
\begin{anmerk}
  Der Binominalkoeffizient ${n\choose k}$ ist demnach ein spezieller 
  Multinominalkoeffizient ${n \choose k_1, k_2}$ mit $k_1=k$ und $k_2=n-k$.
\end{anmerk}


\section{Die Rolle der Entropie bei asymptotischen Ab\-sch"atzungen}
Wir f"uhren den Begriff der Entropie ein und kl"aren ihre Beudeutung bei
asympotischen Absch"atzungen.

\subsection{Kompositionsklassen}
Mit ${\cal A}_q$ bezeichnen wir ein endliches \emph{Alphabet}, d.h. eine 
endliche Menge 
$$ {\cal A}_q=\{a_1, a_2,\ldots, a_q\} \, ,$$
die $q$ verschiedene Symbole $a_1, a_2,\ldots, a_q$ enth"alt. Wir betrachten 
Sequenzen der L"ange $n$ "uber ${\cal A}_q$. Die Anzahl der Sequenzen w"achst
exponentiell mit der L"ange $n$.

"Ublicherweise werden die Sequenzen \emph{lexikographisch} geordnet, wobei
$$ a_1 < a_2 < \ldots < a_q $$
gilt. Wir betrachten  jetzt eine andere Ordnung, die sp"ater eine 
zentrale Rolle spielen wird. Dazu definieren wir f"ur jede Sequenz "uber
dem Alphabet ${\cal A}_q$ der L"ange $n$ das $q$-Tupel
$$ \vec{k}=(k_1, k_2, \ldots, k_q) \mbox{ mit } 
  \sum_{i=1}^n k_i=n \, ,$$
wobei $k_i$ angibt, wie oft das Symbol $a_i$ in der Sequenz vorkommt
(die Reihenfolge ist dabei egal). Eine \emph{Kompositionsklasse} besteht
aus Sequenzen, deren $\vec{k}$'s "ubereinstimmen. Die Anordnung der
Sequenzen innerhalb einer Kompositionsklasse nicht entscheidend; wir werden
sie lexikographisch anordnen. Wir ordnen die Kompositionsklassen 
lexikographisch bez"uglich ihrer zugeh"origen $\vec{k}$-Vektoren an.
F"ur gegebene $q$ und $n$ ist eine Kompositionsklasse durch ihr zugeh"origes
$q$-Tupel $\vec{k}$ eindeutig bestimmt.
%
\begin{definition}
  Die Anzahl der verschiedenen Kompositionsklassen f"ur gegebene $q$ und $n$
  wird mit ${\cal K}(q,n)$ bezeichnet.
\end{definition}
%
\begin{satz}
Die Anzahl der verschiednen Kompositionsklassen ${\cal K}(q,n)$ ist die 
Anzahl der verschiedenen Kombinationen mit Wiederholung von $n$ Elementen der 
Ordnung $q$ und damit 
$$ {\cal K}(q,n)=A(\tilde{C}_q^n) = {n+q-1\choose q} \, . $$
\end{satz}
%
\begin{definition}
Die Anzahl der Sequenzen in der Kompositionsklasse mit dem zugek"origen Tupel
hei"st \emph{Multiplizit"at} $M(\vec{k})$ der Kompositionsklasse.
\end{definition}
%
\begin{satz}
Die Multiplizit"at $M(\vec{k})$ der Kompositionsklasse mit dem zugeh"origen
Tupel $\vec{k}=(k_1,k_2,\ldots,k_q)$ ist durch
den Multinominalkoeffizienten 
$$ M(\vec{k}) = {n \choose k_1, k_2,\ldots, k_q}\,, $$
gegeben, wobei $n=\sum_{i=1}^q k_i$.
\end{satz}


\subsection{Herleitung der bin"aren Entropiefunktion}
Wir betrachten hier den Fall $q=2$ und bezeichnen die beiden Symbole mit
$A$ und $B$. Wir werden die bin"are Entropiefunktion herleiten und ihre 
Bedeutung bei asymptotischen Absch"atzungen erkl"aren. Diese Ergebnisse werden 
im n"achsten Abschnitt f"ur $q>2$ verallgemeinert.

Wir wollen das asymptotische Verhalten der Gr"ossen von Kompositionsklassen 
f"ur $n\rightarrow\infty$ untersuchen, wobei die relative H"aufigkeit der 
$A$'s gleich bleiben soll. Wir betrachten also eine Folge von 
Kompositionsklassen mit den Paaren $\vec{k}=(k_A,k_B)$, wobei 
$k_A=\lambda n$ und $k_B = (1-\lambda) n$. Wir wissen, da"s die Gr"ossen 
durch die Binominalkoeffizienten ${n\choose \lambda n}$ gegeben sind. Wir
benutzten die Approximationen 
$$
  \sqrt{2\pi n} \cdot n^n \cdot e^{-n} \cdot e^{\frac{1}{12n+1}} < n ! <
  \sqrt{2\pi n} \cdot n^n \cdot e^{-n} \cdot e^{\frac{1}{12n}} \, ,
$$
um Absch"atzungen f"ur die Binominalkoeffizienten zu bekommen.

Als untere Schranke f"ur den Binomialkoeffizienten 
$$ {n \choose k} = {n! \over k! (n-k)!} $$
erhalten wir
$$ {1\over\sqrt{2\pi n \lambda (1-\lambda)}} \cdot
   n^n \cdot k^{-k} \cdot (n-k)^{-(n-k)} \cdot
   e^{{1\over 12n+1} - {1\over 12k} - {1\over 12(n-k)}}, $$
indem wir f"ur den Z"ahler $n!$ die obige Unterschranke und f"ur die Faktoren
$k!$ und $(n-k)!$ im Nenner die Oberschranke verwenden.
Der Logarithmus des Terms $n^n \cdot k^{-k} \cdot (n-k)^{-(n-k)}$ ist
$$ n\log_2 n - k\log_2 k - (n-k)\log_2 (n-k).$$
Wenn wir $n\log_2 n$ in die Summe $k\log_2 n + (n-k)\log_2 n$ aufspalten und
die Logarithmenregeln anwenden, vereinfacht sich der Ausdruck zu
$$ k\log {n\over k} + (n-k)\log_2 (n-k). $$
Wir verwenden jetzt, da"s die relative H"aufigkeit $\underline{k}_A=\lambda$ 
fest bleibt, d.h.
$k=k_A=\lambda n$, und damit
$$ n\left
      (\lambda\log_2{1\over\lambda} + (1-\lambda)\log_2{1\over 1-\lambda}
    \right) =
   n H_2(\lambda), $$
wobei
$$
  H_2(\lambda) = -\lambda \log_2 \lambda - (1-\lambda) \log_2 (1-\lambda)
$$
die \emph{bin"are Entropiefunktion} bezeichnet. Die \emph{Unterschranke} f"ur
den Binominalkoeffizienten ist also
$$
  {1\over\sqrt{2\pi n \lambda (1-\lambda)}} \cdot 
  2^{n H_2(\lambda)} \cdot 
  e^{{1\over 12n+1} - {1\over 12\lambda n} - {1\over 12(1-\lambda)n}}
  \le {n \choose \lambda n}.
$$
Die \emph{Oberschranke} wird analog hergeleitet. Sie ist
$$
  {n \choose \lambda n} \le
  {1\over\sqrt{2\pi n \lambda (1-\lambda)}} \cdot 
  2^{n H_2(\lambda)} \cdot
  e^{{1\over 12n} - {1\over 12\lambda n+1} - {1\over 12(1-\lambda)n+1}}.
$$
Wir k"onnen die Oberschranke auch als
$$
  2^{n H_2(\lambda)} \cdot
  2^{-{1\over 2} \log_2 (2\pi n \lambda (1-\lambda)) +
    \log_2 e \cdot ({1\over 12n} - {1\over 12\lambda n+1} - 
    {1\over 12(1-\lambda)n+1})}
$$
schreiben, indem wir den Logarithmus von der Wurzel ziehen und ihn im 
Exponenten schreiben. Der Term 
$$  {1\over 2} \log_2 (2\pi n \lambda (1-\lambda)) -
    \log_2 e \cdot \left({1\over 12n} - {1\over 12\lambda n+1} - 
    {1\over 12(1-\lambda)n+1}\right) $$
kann in der Funktionsklasse $o(n)$ versteckt werden, wobei
$$ o(n) := \left\{ f : \N \rightarrow \R \, , \lim_{n\rightarrow\infty}
                   {f(n)\over n} = 0 \right\} \, . $$
Damit haben wir den folgenden Satz bewiesen:
\begin{satz}[Die bin"are Entropiefunktion]\label{binaereEntropie}
Es gilt:
$$
  {n\choose \lambda n} = 2^{n H_2(\lambda)-o(n)}
$$
Genauer:
$$
  2^{n H_2(\lambda)-f(n)} \le {n\choose \lambda n} 
  \le 2^{n H_2(\lambda)-g(n)} \, ,
$$
wobei 
\begin{eqnarray*}
  f(n) & = & 
  {1\over 2} \log_2 (2\pi n \lambda (1-\lambda)) -
    \log_2 e \cdot \left({1\over 12n+1} - {1\over 12\lambda n} - 
    {1\over 12(1-\lambda)n}\right) \\
  g(n)  & = & 
  {1\over 2} \log_2 (2\pi n \lambda (1-\lambda)) -
    \log_2 e \cdot \left({1\over 12n} - {1\over 12\lambda n+1} - 
    {1\over 12(1-\lambda)n+1}\right)
\end{eqnarray*}
\end{satz}


\subsection{Die $q$-n"are Entropiefunktion}
\begin{definition}[Die $q$-n"are Entropiefunktion]
Es sei
$\vec{\underline{k}}=(\underline{k}_1,\underline{k}_2,\ldots,\underline{k}_q)$
mit $\underline{k}_i\ge 0$ und $\sum_{i=1}^q \underline{k}_i=1$.
Die Funktion
$$
  H_q(\vec{\underline{k}}) = 
  \sum_{i=1}^q \underline{k}_i\log \underline{k}_i
$$
hei"st die \emph{$q$-n"are Entropiefunktion}. Mit dem $q$-Tupel
$\vec{\underline{k}}$ werden die relativen H"aufigkeiten der Alphabetsymbole
$a_1,a_2,\ldots,a_q$ festgelegt.
\end{definition}
"Ahnlich im bin"aren Fall kann die $q$-n"are Entropiefunktion heranzienhen, um
Multiplizit"aten von Kompositionsklassen zu abzusch"aten:
$$ M(\vec{k})={n \choose k_1, k_2, \ldots, k_q} \approx 
2^{n H_q(\vec{\underline{k}})}\,. $$


\section{Entropie als Ungewi"sheit}
Wir haben im vorhergehenden Abschnitt die Entropie hergeleitet und ihre
Bedeutung bei asymptotischen Absch"atzungen erl"autert. Jetzt wollen wir
zeigen, da"s die Entropie auch ein Ma"s f"ur Ungewi"sheit darstellt. Dieser
Zugang wird "ublicherweise in der Informationstheorie benutzt. Betrachten wir
einige Beispiele, die verdeutlichen, was eigentlich Ungewi"sheit bedeutet:
\begin{itemize}
  \item Der Ausgang eines Rennens zwischen zwei ebenb"urtigen Pferden ist
    nicht so ungewi"s wie der eines Rennens zwischen sechs ebenb"urtigen
    Pferden.
  \item Beim Roulette ist das Ergebnis ungewisser als bei einem W"urfelwurf.
  \item Wird ein unverf"alschter W"urfel geworfen, ist das Ergebnis 
    ungewisser, als wenn ein einseitig belasteter W"urfel geworfen wird,
    bei dem die Wahrscheinlichkeit ${1\over 10}$ besteht, eine der Zahlen
    $1$ bis $5$ zu w"urfeln, und ${1\over 2}$, eine $6$ zu w"urfeln.
\end{itemize}
Wir fordern axiomatisch einige Eigenschaften, die man sinnvollerweise von
der Ungewi"sheit erwartet und zeigen dann, da"s die Entropiefunktion die einzige
Funktion ist, die diese Eigenschaften erf"ullt. Deswegen stellt die 
Entropie ein Ma"s f"ur die Ungewi"sheit dar.
\subsection*{Eigenschaften der Ungewissheit}
\begin{itemize}
  \item[E1] $H$ ist eine symmetrische Funktion der Argumente 
    $p_1,p_2,\ldots,p_q$, d.h. 
    f"ur jede Permutation $\pi$ von $1,2,\ldots,q$ gilt:
    $$  H(p_1,p_2,\ldots,p_q)=H(p_{\pi(1)},p_{\pi(1)},\ldots,p_{\pi(1)})\,. $$
    Es kommt lediglich auf die Wahrscheinlichkeiten an und nicht auf deren
    Reihenfolge.
  \item[E2] $H$ ist eine nichtnegative Funktion, d.h. es gilt:
    $$ H(p_1,\ldots,p_q)\ge 0\,. $$
    $H$ nimmt den Wert Null an, genau dann wenn
    $$ \exists i : p_i=1 \,. $$
    Die Ungewissheit ist eine von Natur aus positive Gr"o"se und nimmt nur
    dann den Wert Null an, wenn keine Zuf"alligkeit vorhanden ist.
  \item[E3] $H$ hat ein Maximum, genau dann wenn
    $$ p_1=p_2=\ldots =p_q={1\over q} \, , $$
    d.h.
    $$ \max_{p_1,\ldots,p_q} H\left(p_1,\ldots,p_q\right) = 
       H\left({1\over q},\ldots,{1\over q}\right) \, . $$
    Der Ausgang eines Rennens ist h"ochst ungewi"s, wenn alle Pferde
    ebenb"urtig sind.
  \item[E4] Es gilt:
    $$ H\left({1\over q},\ldots,{1\over q}\right) \le 
       H\left({1\over q+1},\ldots,{1\over q+1}\right) $$
    Der Ausgang eines Rennens mit zwei ebenb"urtigen Pferden ist weniger
    ungewi"s als der eines Rennens mit drei ebenb"urtigen Pferden.
  \item[E5] Es gilt:
    $$ H\left({1\over q_1 q_2},\ldots,{1\over q_1 q_2}\right) =
       H\left({1\over q_1},\ldots,{1\over q_1}\right) + 
       H\left({1\over q_2},\ldots,{1\over q_2}\right) $$
    Die Gesamtungewi"sheit der Ausg"ange von zwei unabh"angigen Pferderennen,
    an denen jeweils ebenb"urtige Pferde teilnehmen, besteht aus der Summe der
    Ungewi"sheiten der einzelnen Pferderennen.
  \item[E6] Es gilt:
    $$ H\left({1\over q},\ldots,{1\over q},0\right) = 
       H\left({1\over q},\ldots,{1\over q}\right)$$
    Der Ausgang eines Rennens, an dem zus"atzlich ein Holzpferd teilnimmt,
    bleibt gleich ungewi"s.
  \item[E7] $H$ sollte eine stetige Funktion ihrer Argumente sein, d.h. bei
    kleinen "Anderungen der Wahrscheinlichkeiten sollte sich die
    Ungewi"sheit nicht wesentlich "andern.
  \item[E8] Es seien $p^{(1)}=p^{(1)}_1+\ldots +p^{(1)}_{q_1}$ und
    $p^{(2)}=p^{(2)}_1+\ldots +p^{(2)}_{q_2}$.
    Dann l"a"st sich bei $p^{(1)}+p^{(2)}=1$ 
    $$ H\left(p^{(1)}_1,\ldots,p^{(1)}_{q_1},
              p^{(2)}_1,\ldots,p^{(2)}_{q_2}\right) $$
    wie folgt zerlegen
    $$ H\left(p^{(1)},p^{(2)}\right) + 
       p^{(1)}\,
H\left({p^{(1)}_1\over p^{(1)}},\ldots ,{p^{(1)}_{q_1}\over p^{(1)}}\right) +
       p^{(2)}\,
H\left({p^{(2)}_1\over p^{(2)}},\ldots ,{p^{(2)}_{q_2}\over p^{(2)}}\right)\,.
    $$
    Die obige Situation tritt beispielsweise bei einem Pferderennen auf, an 
    dem zusammen $q_1$ schwarze Pferde und $q_2$ weisse Pferde teilnehmen.
    Dabei ist $p^{(1)}_i$ die Wahrscheinlichkeit, da"s das $i$-te 
    schwarze Pferd gewinnt, und $p^{(2)}_j$ die Wahrscheinlichkeit, da"s das 
    $j$-te weisse Pferd gewinnt. Folglich sind $p^{(1)}$ und
    $p^{(2)}$ die Wahrscheinlichkeiten, da"s ein schwarzes bzw. ein weisses 
    Pferd gewinnt.

    Die Gesamtungewi"sheit ergibt sich nun aus Summe der Ungewi"sheit, ob ein 
    schwarzes oder weisses Pferd gewinnt, sowie den gewichteten Ungewi"sheiten 
    unter der Voraussetzung, da"s der Gewinner ein schwarzes bzw. weisses Pferd 
    ist.
\end{itemize}
%
Das folgende Theorem zeigt, da"s die Entropiefunktion die einzige Funktion
ist, die die Eigenschaften E1-E8 der Ungewi"sheit erf"ullt.
%
\begin{theorem}[Entropie als Ma"s f"ur Ungewi"sheit]
  Es sei $H(p_1,\ldots,p_q)$ eine Funktion, die f"ur alle Zahlen
  $p_1,\ldots,p_q\in \mathbb{R}$ mit
  $$ p_i\ge 0 \mbox{ und } \sum_{i=1}^q p_i = 1 $$
  definiert ist. Wenn $H$ die Eigenschaften E1 bis E8 erf"ullen soll, dann gilt
  $$ H(p_1,\ldots,p_q)=-\lambda\sum_{i=1}^q p_i\log_b p_i \, .$$
  Hierbei ist $\lambda\in\mathbb{R}$ eine beliebige positive Konstante, und die
  Summe wird "uber diejenigen $i$ gebildet, f"ur die $p_i>0$ gilt.
\end{theorem}
Beweis: Wir definieren
$$ g(q)=H\left({1\over q},{1\over q},\ldots,{1\over q}\right) $$
f"ur alle $q\in\mathbb{N}$. 
Nach E5 
$$ H\left({1\over q_1 q_2},\ldots,{1\over q_1 q_2}\right) =
   H\left({1\over q_1},\ldots,{1\over q_1}\right) + 
   H\left({1\over q_2},\ldots,{1\over q_2}\right) $$
gilt (setze $q_1=q_2=q$)
$$ g(q^2)=g(q)+g(q)=2 g(q) $$
und damit
$$ g(q^k)=k\cdot g(q) $$
f"ur $k\in\mathbb{N}$. Nach E4
$$ H\left({1\over q},\ldots,{1\over q}\right) \le
   H\left({1\over q+1},\ldots,{1\over q+1}\right) $$
gilt 
$$ g(q)\le g(q+1) $$  
und damit sieht man, da"s $g$ eine monoton steigende Funktion, d.h.
$$ g(q_1)\le g(q_2)$$
f"ur $q_1\le q_2$ und $q_1,q_2\in\mathbb{N}$. Funktionen, die 
diese beiden Eigenschaften erf"ullen, sind die elementaren Funktionen
$$ g(q)=A\log_b(q) \, $$
wobei $A\in\R^+$ eine Konstante ist und der Logarithmus zu einer beliebigen
Basis $b>1$ genommen wird.

Als n"achstes betrachten wir die Eigenschaft E8: F"ur 
$p^{(1)}=p^{(1)}_1+\ldots +p^{(1)}_{q_1}$
und
$p^{(2)}=p^{(2)}_1+\ldots +p^{(2)}_{q_2}$ l"a"st sich 
$$ H\left(p^{(1)}_1,\ldots,p^{(1)}_{q_1},p^{(2)}_1,\ldots,p^{(2)}_{q_2}\right)
$$
in
$$ H\left(p^{(1)},p^{(2)}\right) + 
   p^{(1)}\cdot 
   H\left({p^{(1)}_1\over p^{(1)}},\ldots,{p^{(1)}_{q_1}\over p^{(1)}}\right) +
   p^{(2)}\cdot
   H\left({p^{(2)}_1\over p^{(2)}},\ldots,{p^{(2)}_{q_2}\over p^{(2)}}\right)
$$
zerlegen. F"ur $q=q_1+q_2$ und 
$$ p^{(1)}_1 = \ldots = p^{(1)}_{q_1} = {1\over q} $$
$$ p^{(2)}_1 = \ldots = p^{(2)}_{q_2} = {1\over q} $$
haben wir
$p^{(1)}={q_1\over q}$ und $p^{(2)}={q_2\over q} = {q-q_1\over q}$. Nach E8
l"ast sich $g(q)$ als 
$$ H\left({q_1\over q},{q-q_1\over q}\right) +
   {q_1\over q} 
   \underbrace{H\left({1\over q_1},\ldots,{1\over q_1}\right)}_{g(q_1)} +
   {q-q_1\over q} 
   \underbrace{H\left({1\over q-q_1},\ldots,{1\over q-q_2}\right)}_{g(q-q_1)}
$$
schreiben. Wir haben also
$$ g(q)=H\left({q_1\over q},{q-q_1\over q}\right)+
   {q_1\over q} g(q_1) + 
   {q-q_1\over q} g(q-q_1)\,. $$
Wenn wir jetzt $g(x)=A\log x$ ausnutzen, so erhalten wir mit elementaren
Umformungen
\begin{eqnarray*}
  H\left({q_1\over q},{q-q_1\over q}\right) 
  & = & 
  A\left(\log q - {q_1\over q}\log q_1 - {q-q_1\over q}\log (q-q_1)\right)\\
  & = &
  -A\left(\log q_1^{{q_1\over q}}+\log(q-q_1)^{{q-q_1\over q}}-\log q\right)\\
  & = &
  -A\log {q_1^{{q_1\over q}} \cdot (q-q_1)^{{q-q_1\over q}}\over q}\\
  & = &
  -A\log\left(
       \left({q_1\over q}\right)^{{q_1\over q}}
  \cdot\left({q-q_1\over q}\right)^{{q-q_1\over q}}
  \right) \\
  & = & 
  -A\log\left({q_1\over q}\log_2\left({q_1\over q}\right) +
              {q-q_1\over q}\log_2\left({q-q_1\over q}\right)
  \right) \\
  & = &
  -A\left(p\log p + (1-p)\log(p-1)\right) \\
  & = &
  H(p,1-p)
\end{eqnarray*}
Damit haben wir f"ur $q=2$ die L"osung.
Wir vermuten nun, da"s allgemein 
\begin{equation}
  H(p_1,p_2,\ldots,p_q) = -A \sum_{i=1}^q p_i\log p_i
\end{equation}
gilt. Wir beweisen dies mit vollst"andiger Induktion. Die 
Induktionsverankerung ist die L"osung f"ur $q=2$. Die Induktionsannahme ist,
da"s die obige Gleichung bis zu einem gewissen $q\in\mathbb{N}$ gilt. Wir
machen den Induktionsschlu"s auf $q+1$ und betrachten dazu
$$ H(p_1,p_2,\ldots,p_q,p_{q+1}) \, . $$
Wir setzen $p_i^{(1)}=p_i$ f"ur $i=1,\ldots,q$, 
$p^{(1)}=\sum_{i=1}^q p_i^{(1)}$ und $p^{(2)}=p_1^{(2)}=p_{q+1}$ und wenden
wieder E8 an. Mit elementaren Umformungen erhalten wir dann
\begin{eqnarray*}
  & &
  H(p_1,\ldots,p_q,p_{q+1}) \\
  & = & 
  H(p^{(1)},p^{(2)}) +
  p^{(1)} \cdot
  H\left({p_1^{(1)}\over p^{(1)}},\ldots,{p_q^{(1)}\over p^{(1)}}\right) +
  p^{(2)} \cdot 
  H\left({p_1^{(2)}\over p^{(2)}}\right) \\
  & = &
  H(p^{(1)},p^{(2)}) +
  p^{(1)} \cdot
  H\left({p_1\over p^{(1)}},\ldots,{p_q\over p^{(1)}}\right) +
  p^{(2)} \cdot 
  \underbrace{H(1)}_{=0} \\
  & = &
  -A \left( p^{(1)}\log p^{(1)} + p^{(2)}\log p^{(2)} \right) -
  p^{(1)}\, A\,\sum_{i=1}^q {p_i\over p^{(1)}} \log {p_i\over p^{(1)}} \\
  & = &
  -A \left( p^{(1)}\log p^{(1)} + p^{(2)}\log p^{(2)} \right) -
  A\,\sum_{i=1}^q p_i \log p_i + A \, p^{(1)} \log p^{(1)} \\
  & = &
  -A \sum_{i=1}^q p_i \log p_i
\end{eqnarray*}
Damit ist die Behauptung f"ur alle $q\in\mathbb{N}$ bewiesen.


\section{Entropie als Informationsma"s}
Es sei $\Omega=\{\omega_1,\ldots,\omega_n\} $ die Menge der Elementarereignisse
und ${\cal P}=(p_1,\ldots,p_q)$ ein Wahrscheinlichkeitsma"s auf $\Omega$.
Die zugeh"orige Entropie ist
$$ H({\cal P}) = H(p_1,\ldots,p_q) = -\sum_{i=1}^q p_i \log p_i \quad
   \left[ {\mbox{bit}\over\mbox{EE}}\right]\,. $$
Wir wollen zeigen, da"s die Entropie der Mittelwert 
der Selbstinformation der Elementarereignisse ist. Dazu
definieren wir ein Ma"s f"ur die Information, die ein
Elementarereignis tr"agt. Betrachten wir die Zufallsvariable 
$$ X=-\log p\,. $$
Zu jedem Elementarereignis $\omega_i$ geh"ort der Wert $x_i$ der
Zufallsvariablen
$$ x_k = -\log P(\omega_i) = -\log p_i\,. $$
Die Gr"o"se $-\log p_i$ wird oft als \emph{Selbstinformation}, das zu dem
Elementarereignis $\omega_i$ geh"ort, bezeichnet:
$$ I(\omega_i)=-\log p_i\,.$$
Die Einheit der Information wird als \emph{Bit} bezeichent, wobei 1 Bit der
der Information entspricht, die mit dem  Ausw"ahlen von zwei 
gleichwahrscheinlichen Elementarereignissen ($p_1,p_2=\frac{1}{2}$) 
assoziiert ist. In diesem Fall haben wir
$$ I(\omega_1)=I(\omega_2)=-\log \frac{1}{2} = 1 \mbox{ bit}\,. $$
Die Auswahl zwischen zwei gleichwahrscheinlichen Elementarereignissen 
ben"otigt 1 bit an Information. Wenn $\Omega$ in $2^n$ gleichwahrscheinliche
Ereignisse partitioniert wird, dann betr"agt die Selbstinformation eines
Elementarereignisses
$$ I(\omega_i) = -\log p_i = -\log 2^{-n} = n \mbox{ bits} $$
Die Selbstinformation ist eine nichtnegative Gr"o"se:
$$ I(\omega_i) = -\log p_i \ge 0\,. $$
Gleichheit ergibt sich nur bei einem sicheren Ereignis; in dem Wissen des
Auftreten dieses Ereignisses liegt kein Informationsgehalt.

Wenn $I(\omega_i)$ den Informationsgehalt angibt, der mit dem Auftreten des
Elementarereignisses verbunden ist, so gibt die Entropie
$$ H(p_1, p_2, \ldots, p_n)=-\sum_{i=1}^n p_i \log p_i = 
   \overline{I(\omega_i)} $$
den mitteren Informationsgehalt an.
\begin{thebibliography}{10}
\bibitem{ham}
  R.~W.~Hamming,
  \emph{Coding and information theory},
  Englewood Cliffs, NJ : Prentice-Hall, 1980.
\bibitem{krengel91}
  U.~Krengel,
  \emph{Einf"uhrung in die Wahrscheinlichkeitstheorie und Statistik},
  vieweg studium,
  1991.
\end{thebibliography}

\end{document}